\chapter{結言}\label{ch:conclusions}
\section{本研究で得られた結論}
本研究では，プール沸騰音響からの熱流束回帰に基づくONB近傍の検知を対象に，異なる入力表現を用いるモデル統合の効果と限界を検討した．2025年6月11日・6月18日の計36 WAVを用い，各録音内の45学習区間・15評価区間という条件で，無雑音学習から7段階の付加雑音条件への固定適用を比較した．以下の結論は，既知WAV内の未使用540区間，1秒，3\,kHz上限，固定基準水流音，seed 42の評価範囲に基づく．

\begin{enumerate}
\item \textbf{異なる要約表現を用いるモデルの統合は，全域回帰と総見逃しの一部を改善した．}
5モデル統合は元3 MSEに対して全7条件で全域RMSEが小さく，無雑音で34.07から27.80\,\kw{}，基準SNR $-20$\,dBで81.14から75.23\,\kw{}となった．総FNもそれぞれ42から27，65から52へ減少した．一方，無雑音の元3＋ETや$-20$\,dBのHGB単体より常に良い結果ではなかった．

\item \textbf{雑音時の絶対性能の改善は，悪化量の縮小と同じではなかった．}
元3 MSEと比べると，無雑音から$-20$\,dBへのRMSE増加量は47.07から47.43\,\kw{}となった．この比較では，雑音時の誤差低減は確認できるが，雑音による悪化量を抑えたとは結論できない．

\item \textbf{全域回帰の改善は，ONB発生段階の十分な検知へは結びつかなかった．}
5モデルの近傍RMSEは無雑音で改善する一方，$-20$\,dBでは元3 MSEより悪化した．記録上ONB段階の30区間について，無雑音で8区間，$-20$\,dBで0区間しか陽性とならなかった．日別$q_{100}$も全方式・両日で一律に前進していない．従って，ONB近傍検知という目的には部分的な利得と未達が併存する．

\item \textbf{統合の利得には，相殺に加えて同方向誤差の縮小が含まれた．}
保存予測の分解により，元3・HGB・ETが同方向に外れた区間群が，$-20$\,dBの全体MSE減少の87.15\%を占めた．誤差の小さい予測へ調整する効果も大きく，逆符号の相殺だけでは補完性を説明できなかった．この割合は予測上の損失分解であり，物理音源の寄与率ではない．

\item \textbf{モデル群の入力依存には差があるが，物理音源の同定は未達である．}
共通の時間行並べ替えでCNN系の誤差が増え，順序を集約するHGB/ETは不変となった．一方，追加2モデルは全域gainや対象帯域の減衰に敏感であり，全ての入力変化に頑健ではなかった．入力依存差と統合の誤差変化は示せるが，水流雑音時の利得の因果割合や，利用帯域の音源は現在の結果だけから確定できない．
\end{enumerate}

以上より，本研究は，モデル統合の効果を全域精度だけで判断せず，ONB判定，熱流束領域，誤差の方向・大きさ，入力依存へ分けることで，改善した点と検知上の限界を具体的に示した．当初目指した早期検知と物理的説明を全面的に達成したものではないが，その達成を判断するための比較と，現在の方法で不足する点を明確にした．

\section{今後の展望}
最初に整える必要があるのは，データ取得者，原計測条件，熱流束ラベルとONBの原実験上の根拠である．これらは新しいモデルを導入しても補えず，解析の再現性と研究上の立場を定めるための確認事項となる．文献については，音響回帰・状態判定の先行研究の本文と，手塚の最終修論における物理的考察を，自研究の原条件と照合して補う．

検知性能をさらに改善する場合には，総FN，FP，開始段階，$q_{100}$のどれを改善対象とするかを先に定める必要がある．全域MSEとは異なる重み決定や入力に応じた自動配分は候補となるが，現結果から自動的な有効性を仮定しない．誤報との交換や，単体の良い判定を統合で失う事例も含め，同じ評価条件で検証する必要がある．

未知日・未知録音への適用を主張する場合は，未参照の独立データと，実運用で利用可能な閾値・雑音基準を定めた評価が必要となる．また，秒単位の検知遅延には気泡発生時刻に対応する正解が，物理音源の確定には同期画像や条件変更による識別根拠が必要である．これらは本稿の限定した結論を拡張するための条件であり，現在の執筆を停止して全て実行するという意味ではない．不足を主張の強さと対応づけ，中心結論に必要な確認から進める．
